%Master File:lectures.tex

\lesson{Sequencing}
\vspace{-2cm}
\begin{center}
 \includegraphics[height=10cm]{cost_per_genome_oct2015.jpg}
\end{center}

\keywords{Shotgun, new sequencing technology, nanopores}
%%%%%%%%%%%%%%%%%%%%%%% Next Slide %%%%%%%%%%%%%%%%%%%%%%%
\renewcommand{\Outline}{%
\begin{slide}
\section[1]{Outline}

\begin{minipage}{12cm}
  \begin{enumerate}\squeeze
    \outlineitem{Sanger Sequencing}{es}
    \outlineitem{Next Generation Sequencers}{newGen}
    \outlineitem{Third Generation Sequencers}{newGen}
   \end{enumerate}
\end{minipage}\hfill
\begin{minipage}{10cm}
  \vfill
  \includegraphics[width=10cm]{cost_per_genome_oct2015.jpg}
  \vfill
\end{minipage}
\end{slide}
\addtocounter{outlineitem}{1}
}

\setcounter{outlineitem}{1}
\Outline
\toptarget{firstoutline}
%%%%%%%%%%%%%%%%%%%%%%% Next Slide %%%%%%%%%%%%%%%%%%%%%%%

\begin{slide}
\section{Sequencing Revolution}

\begin{PauseHighLight}
  \begin{itemize}
  \item Advances in our ability to sequence DNA lies at the heart of
    the biological revolution\pause
  \item Despite the structure of DNA being discovered in 1953, reading
    DNA sequences took several decades to take off\pause
  \item Francis Crick had speculated on how DNA might code for
    proteins\pause
  \item By removing some number of bases from DNA he was able to show
    that DNA has a three letter code\pause
  \end{itemize}
\end{PauseHighLight}

\end{slide}

%%%%%%%%%%%%%%%%%%%%%%% Next Slide %%%%%%%%%%%%%%%%%%%%%%%

\begin{slide}
\section{Frederick Sanger}

\begin{minipage}{0.3\linewidth}
  \includegraphics[width=0.9\linewidth]{Frederick_Sanger2.jpg}
\end{minipage}\hfil
\begin{minipage}{0.67\linewidth}
  \begin{PauseHighLight}
  \begin{itemize}
  \item Frederick Sanger was a British born biochemist\pause
  \item In 1955 he discovered the protein sequence for Insulin\pause
  \item This showed that proteins were actually chains\pause
  \item He won the Nobel prize for chemistry for this in 1958\pause
  \end{itemize}
\end{PauseHighLight}

\end{minipage}
\end{slide}

%%%%%%%%%%%%%%%%%%%%%%% Next Slide %%%%%%%%%%%%%%%%%%%%%%%

\begin{slide}
\section{Sequencing DNA}

\begin{PauseHighLight}
  \begin{itemize}
  \item Chemically it was difficult because the bases are not that
    distinctive\pause
  \item It wasn't until 1965 that the very first strands of RNA were
    read\pause
  \item A lot of effort went into learning how to copy DNA and break it
    up into smaller sequences\pause
  \item Various techniques were developed in the 1970's to sequence
    longer lengths of DNA until in 1977 Sanger developed a new method
    now known as \emph{``Sanger Sequcencing''}\pause
  \end{itemize}
\end{PauseHighLight}

\end{slide}

%%%%%%%%%%%%%%%%%%%%%%% Next Slide %%%%%%%%%%%%%%%%%%%%%%%

\begin{slide}
\section{Sanger Sequencing}

\begin{PauseHighLight}
  \begin{itemize}
  \item In Sanger sequencing we use biological mechanisms for copying
    DNA\pause
  \item This would add Dideoxynucleoties (ddNTP) to a DNA chain\pause
  \end{itemize}
  \begin{center}
    \includegraphics[width=0.6\linewidth]{DNA_Replication}
  \end{center}
\end{PauseHighLight}

\end{slide}


%%%%%%%%%%%%%%%%%%%%%%% Next Slide %%%%%%%%%%%%%%%%%%%%%%%

\begin{slide}
\section{Modified ddNTP}

\begin{PauseHighLight}
  \begin{itemize}
  \item Modified ddNTP's were developed with a missing $3'$ hydroxyl
    group which prevented the chain from developing any further\pause
  \item These were marked (either with radioactive phosphorous or later
    with a fluorescent marker)\pause
  \item DNA replication was then carried out with normal ddNTP's and
     a weak concentration of these modified ddNTP's\pause
  \item Initially this was done in four containers one for each base
    (in container A only the base $A$ ddNTP were modified)\pause
  \item Later different fluorescence were used so everything could be
    carried out in the same vessel\pause
  \end{itemize}
\end{PauseHighLight}

\end{slide}

%%%%%%%%%%%%%%%%%%%%%%% Next Slide %%%%%%%%%%%%%%%%%%%%%%%

\begin{slide}
\section{Sanger Sequencing}
\pause
\pb
\pauselevel{=1}
\begin{center}
  \multipdf[width=\linewidth]{sangerMethod}\pause
\end{center}
\end{slide}

%%%%%%%%%%%%%%%%%%%%%%% Next Slide %%%%%%%%%%%%%%%%%%%%%%%

\begin{slide}
\section[-2]{Fractionation}

\begin{minipage}{0.3\linewidth}
  \includegraphics[width=\linewidth]{sangerMethod1}
\end{minipage}\hfil
\begin{minipage}{0.65\linewidth}
  \begin{PauseHighLight}
  \begin{itemize}
  \item Each solution is added to a layer of gel\pause
  \item Using capillary electrophoresis different length DNA fragments
    are sorted\pause
  \item From the distance travelled and the colour of the last
    nucleotide base we can read off the sequence\pause
  \item This won Frederick Sanger a second Nobel prize awarded in
    1980\pause
  \end{itemize}
\end{PauseHighLight}

\end{minipage}
\end{slide}

%%%%%%%%%%%%%%%%%%%%%%% Next Slide %%%%%%%%%%%%%%%%%%%%%%%

\begin{slide}
\section[-1]{Shotgun Sequencing}

To save careful preparation researcher developed shotgun sequencing
\begin{center}
  \includegraphics[width=0.9\linewidth]{shotgun}
\end{center}
\end{slide}

%%%%%%%%%%%%%%%%%%%%%%% Next Slide %%%%%%%%%%%%%%%%%%%%%%%

\begin{slide}
\section{Reconstruction}

\begin{PauseHighLight}
  \begin{itemize}
  \item Reconstruction looks easy, but the human genome is full of very
    repetitive sequence\pause
  \item Sanger sequences only allowed you to sequence ``reads'' of
    around 1000bp\pause
  \item Not clear what the long range ordering is\pause
  \item For large genomes researchers built ``scaffolds'' where they
    knew the position of some genes\pause
  \item You can sequence both sides of a large fragment so you know the
    approximate relationship between reads\pause
  \end{itemize}
\end{PauseHighLight}

\end{slide}

%%%%%%%%%%%%%%%%%%%%%%% Next Slide %%%%%%%%%%%%%%%%%%%%%%%

\begin{slide}
\section{Human Genome Project}

\begin{PauseHighLight}
  \begin{itemize}
  \item Sanger sequencing was the technology that delivered the human
    genome\pause
  \item Huge advances meant that it was delivered several years earlier
    than expect\pause
  \item The cost of sequencing a human genome in 2003 was estimated at
    around \$10-100 million\pause
  \item The human genome project did not deliver an immediate benefit to
    medicine that some had anticipated---we still did not and don't
    understand enough biology\pause
  \end{itemize}
\end{PauseHighLight}

\end{slide}




%%%%%%%%%%%%%%%%%%%%%%% Next Slide %%%%%%%%%%%%%%%%%%%%%%%
\Outline % Next Generation
%%%%%%%%%%%%%%%%%%%%%%% Next Slide %%%%%%%%%%%%%%%%%%%%%%%

\begin{slide}
\section{Improving on Sanger}

\begin{PauseHighLight}
  \begin{itemize}
  \item Although Sanger sequencing saw improves in efficiency in line
    with Moore's law there was a desire to reduce the cost of sequencing
    a human genome to \$1000 within a decade\pause
  \item The ideas was to use massive parallelisation\pause
  \item Many different techniques were tried\pause
  \end{itemize}
\end{PauseHighLight}

\end{slide}

%%%%%%%%%%%%%%%%%%%%%%% Next Slide %%%%%%%%%%%%%%%%%%%%%%%

\begin{slide}
\section[-2]{Sequencing Techniques}

\begin{PauseHighLight}
  \begin{itemize}
  \item Sanger sequencing used sequencing-by-synthesis\pause
    \begin{itemize}
    \item Using DNA-polymerase a strand of DNA is replicated\pause
    \item It turns out this technique that ultimately triumphant\pause
    \end{itemize}
  \item Another approach is sequencing-by-ligation (or hybridisation)\pause
    \begin{itemize}
    \item Create lots of short sections of DNA (oligonucleotides) that
      will stick to an unbound stretch of DNA (or RNA)\pause
    \item Used in microarrays\pause
    \end{itemize}
  \end{itemize}
\end{PauseHighLight}

\end{slide}

%%%%%%%%%%%%%%%%%%%%%%% Next Slide %%%%%%%%%%%%%%%%%%%%%%%

\begin{slide}
\section[-2]{Synthesis versus Ligation}

\begin{center}
  \includegraphics[width=0.8\linewidth]{ligationvsynthesis}
\end{center}
\end{slide}


%%%%%%%%%%%%%%%%%%%%%%% Next Slide %%%%%%%%%%%%%%%%%%%%%%%

\begin{slide}
\section{454 Sequencer}

\begin{PauseHighLight}
  \begin{itemize}
  \item The first high throughput sequencer was the Roche 454
    machine\pause
  \item It amplified fragments of DNA attached to a microbead\pause
  \item These are contained in micro-wells\pause
  \item Sequencing-by-synthesis is carried out\pause
  \item Added one type of base at a time with fluorescent
    label\pause
  \item Can have difficulty distinguishing regions with repeated base\pause
  \end{itemize}
\end{PauseHighLight}

\end{slide}

%%%%%%%%%%%%%%%%%%%%%%% Next Slide %%%%%%%%%%%%%%%%%%%%%%%

\begin{slide}
\section[-2]{Water in Oil Emulsion PCR}

\begin{center}
  \includegraphics[width=0.9\linewidth]{454}
\end{center}
\end{slide}


%%%%%%%%%%%%%%%%%%%%%%% Next Slide %%%%%%%%%%%%%%%%%%%%%%%

\begin{slide}
\section[-2]{Illumina}

\begin{PauseHighLight}
  \begin{itemize}
  \item The dominant next generation machines are produced by Illumina
    (formally Solexa)\pause
  \item It uses a technique now called bridge amplification to create
    clusters of identical strands on a 2-D surface (substrate)\pause
  \item The machine can sequence 100's of million clusters at once\pause
  \item It then uses sequencing-by-synthesis to read the sequence\pause
  \item It can also read the other end of the sequence by flipping the
    sequence\pause
  \item Try Googling ``Illumina sequencing videos''\pause
  \end{itemize}
\end{PauseHighLight}

\end{slide}

%%%%%%%%%%%%%%%%%%%%%%% Next Slide %%%%%%%%%%%%%%%%%%%%%%%

\begin{slide}
\section[-1]{Bridge Amplification}

\begin{center}
  \includegraphics[width=\linewidth]{bridging-pcr-large}
\end{center}
\end{slide}

%%%%%%%%%%%%%%%%%%%%%%% Next Slide %%%%%%%%%%%%%%%%%%%%%%%

\begin{slide}
\section[-2]{Identifying Sequence}

\begin{center}
  \includegraphics[width=0.9\linewidth]{illumina-synthesis}
\end{center}
\end{slide}


%%%%%%%%%%%%%%%%%%%%%%% Next Slide %%%%%%%%%%%%%%%%%%%%%%%

\begin{slide}
\section[-2]{Cost Per Human Genome}

\begin{center}
 \includegraphics[width=0.9\linewidth]{cost_per_genome_oct2015.jpg}
\end{center}

\end{slide}


%%%%%%%%%%%%%%%%%%%%%%% Next Slide %%%%%%%%%%%%%%%%%%%%%%%

\begin{slide}
\section{100\,000 Genome Project}

\begin{PauseHighLight}
  \begin{itemize}
  \item Currently the NHS and Sanger centre is sequencing around 70\,000
    human genomes to identify rare genetic diseases\pause
  \item We can now sequence a human genome within a day\pause
  \item Because we have already sequenced the human genome we no longer
    have difficulty reconstructing the sequence\pause
  \end{itemize}
\end{PauseHighLight}

\end{slide}



%%%%%%%%%%%%%%%%%%%%%%% Next Slide %%%%%%%%%%%%%%%%%%%%%%%
\Outline % Third Generation
%%%%%%%%%%%%%%%%%%%%%%% Next Slide %%%%%%%%%%%%%%%%%%%%%%%

\begin{slide}
\section{Single Molecule DNA Sequencing}

\begin{PauseHighLight}
  \begin{itemize}
  \item Having to replicate DNA can bring in errors and biases\pause
  \item Most of the high-throughput techniques only read short sections
    ($\approx100$ base pairs), which can make reconstruction
    difficult\pause
  \item The current research is to create a machine that can read a
    long single DNA strand\pause
  \end{itemize}
\end{PauseHighLight}

\end{slide}

%%%%%%%%%%%%%%%%%%%%%%% Next Slide %%%%%%%%%%%%%%%%%%%%%%%

\begin{slide}
\section{Zero Mode Waveguide}

\begin{minipage}{0.5\linewidth}
  \includegraphics[width=\linewidth]{singleMolecule}
\end{minipage}\hfill
\begin{minipage}{0.47\linewidth}
  \begin{PauseHighLight}
  \begin{itemize}
  \item Uses same fluorescence dNTPs as used in high-throughput
    processing\pause
  \item Use micro-well smaller than the wavelength of light\pause
  \item Able to detect very small signal\pause
  \item Not high-throughput\pause
  \end{itemize}
\end{PauseHighLight}

\end{minipage}
\end{slide}

%%%%%%%%%%%%%%%%%%%%%%% Next Slide %%%%%%%%%%%%%%%%%%%%%%%

\begin{slide}
\section{Nanopore Technology}

\begin{PauseHighLight}
  \begin{itemize}
  \item One of the most anticipated new technologies is to force DNA
    through a tiny pore and read out the sequence\pause
  \item Oxford nanopores have been at the forefront of
    developments\pause
  \item They have just released a hand-held machine\pause
  \item Currently the accuracy is not quite good enough\pause
  \end{itemize}
\end{PauseHighLight}

\end{slide}


%%%%%%%%%%%%%%%%%%%%%%% Next Slide %%%%%%%%%%%%%%%%%%%%%%%

\begin{slide}
\section[-2]{Nanopore}

\begin{center}
  \includegraphics[width=0.9\linewidth]{nanopore1}
\end{center}
\end{slide}

%%%%%%%%%%%%%%%%%%%%%%% Next Slide %%%%%%%%%%%%%%%%%%%%%%%

\begin{slide}
\section{Nanopores}

\begin{center}
  \includegraphics[width=0.9\linewidth]{nanopore2}
\end{center}
\end{slide}

%%%%%%%%%%%%%%%%%%%%%%% Next Slide %%%%%%%%%%%%%%%%%%%%%%%

\begin{slide}
\section{Conclusion}

\begin{PauseHighLight}
  \begin{itemize}
  \item One of the dramatic revolutions that has occurred in your life
    times has been the ability to sequence DNA\pause
  \item Although sequencing proteins predates the sequencing of DNA
    their has been no such revolution\pause
  \item The technology builds on the natural tools of the cell (e.g. DNA
    replicase)\pause
  \item Computationally the ability to reassemble reads into longer
    ``contigs'' has been essential to this revolution\pause
  \end{itemize}
\end{PauseHighLight}

\end{slide}



